\subsection{Separate clinical targeted-assay stress test}
We retain the same 30 genes in MSK-IMPACT 2017 \cite{zehir}, a mixed-cancer clinical cohort with 10,945 specimens from 10,336 patients, including 562 patients with repeated specimens and up to five specimens per patient. The separate addendum was locally recorded before complete outcome retrieval. DataHub sources are pinned to commit \texttt{dca75cb}; downloaded large-file content hashes and byte counts match its Git LFS pointers. A gene is profiled only if the specimen's declared mutation panel contains it. Primary context uses the source \texttt{SAMPLE\_TYPE=Primary}; one unspecified type remains in all-specimen analyses and is not imputed primary.

Both counting engines agree on all 480 specifications. All selected genes occur in both IMPACT341 and IMPACT410, so this panel has no gene-unprofiled specimens: observed and completion certificates coincide, and every fixed-specification point interval is a singleton. All 30 families are sensitive; reference class flips are sensitive in 30/30, context flips in 8/30, and unit/assay flips in 0/30 (Fig.~\ref{fig:msk}). The largest reference unit effect is 0.552 pp and 4.924\% relative. The optimal, edge-greedy, best global fixed, and best per-family fixed policies all average 1.500 questions; entropy-only splitting requires 4.000. This tests repeated-specimen counting and a targeted-panel adapter, but supplies no additional empirical missing-status or adaptive-planning advantage. It is a distinct clinical cohort in the same curation ecosystem, not an independent portal or a pooled TCGA prevalence estimate.

